\section*{Chapitre \ref{ch:m3:on-chain-order-books-and-perpetual-dexs} --- Carnets d'ordres on-chain et DEX de perpétuels}

\begin{solution}{exo:m3:on-chain-order-books-and-perpetual-dexs:1}
Ordre d'arrivée~: l'IOC achète le cours vendeur à 101, l'annulation ne trouve rien, et le nouveau cours vendeur reste au carnet à 103. \omterm{def:m3:on-chain-order-books-and-perpetual-dexs:priority}{Règle de priorité dans le bloc}~: le nouveau cours vendeur à 103 reste, l'annulation
retire le cours vendeur à 101, et l'IOC ne trouve aucun cours vendeur à 101 ou en dessous et expire. Seule la cotation du teneur de marché à 103 subsiste dans les deux cas.
\end{solution}

\begin{solution}{exo:m3:on-chain-order-books-and-perpetual-dexs:2}
Chaque cotation et chaque annulation coûteraient du \omterm{def:m3:blockchains-for-traders:gas}{gas} et attendraient jusqu'à un bloc~: les teneurs de marché ne pourraient ni se permettre de mettre à jour leurs cotations ni annuler avant de subir
une exécution opportuniste~; ils coteraient donc large, ou pas du tout~; et le \omterm{def:m3:blockchains-for-traders:mempool}{mempool} public montrerait chaque annulation avant son exécution.
\end{solution}

\begin{solution}{exo:m3:on-chain-order-books-and-perpetual-dexs:3}
Un \omterm{def:m3:blockchains-for-traders:pow}{validateur}, ou une minorité d'enjeu, publiant un faux prix~: la médiane ignore les valeurs aberrantes, et la déplacer exige qu'une majorité de l'enjeu bouge ensemble. La médiane pondérée
calculée sur les plateformes retenues par chaque \omterm{def:m3:blockchains-for-traders:pow}{validateur} pour son prix, protège contre le déplacement du prix d'une seule plateforme.
\end{solution}

\begin{solution}{exo:m3:on-chain-order-books-and-perpetual-dexs:4}
La marge initiale à 20 fois est de 5~\%, donc la marge de maintenance est de 2.5~\%~; le \omterm{def:m3:on-chain-order-books-and-perpetual-dexs:vault}{coffre de liquidité} prend le relais sous les deux tiers de celle-ci, soit des fonds propres d'environ 1.67~\% du notionnel.
\end{solution}

\begin{solution}{exo:m3:on-chain-order-books-and-perpetual-dexs:5}
$4\text{ millions} \times (1 - 0.111) = $ 3.56~millions de dollars.
\end{solution}

\begin{solution}{exo:m3:on-chain-order-books-and-perpetual-dexs:6}
Dans un \omterm{def:m3:automated-market-makers:amm}{pool de liquidité} à prix d'\omterm{def:m3:blockchains-for-traders:contract}{oracle}, le fournisseur de liquidité est la contrepartie de chaque transaction au prix de l'\omterm{def:m3:blockchains-for-traders:contract}{oracle}~; il perd face à quiconque connaît le prix mieux ou plus tôt
(latence de l'\omterm{def:m3:blockchains-for-traders:contract}{oracle}, manipulation), auxquelles s'ajoute le P\&L net des traders. Dans un coffre, il supporte l'inventaire de son activité de tenue de marché et la queue de risque des liquidations reprises.
Ces deux structures concentrent le risque de la plateforme sur les déposants~; celui du pool est continu, celui du coffre arrive par sauts.
\end{solution}

\begin{solution}{exo:m3:on-chain-order-books-and-perpetual-dexs:7}
Ordre d'arrivée~: 980 lots exécutés de façon opportuniste et 1\,480 pas de cotation perdus~; \omterm{def:m3:on-chain-order-books-and-perpetual-dexs:priority}{règle de priorité dans le bloc}~: aucun. La protection de la règle ne dépend pas de celui qui arrive le premier
dans le bloc, mais seulement de la présence de l'annulation du teneur de marché dans ce même bloc.
\end{solution}

\begin{solution}{exo:m3:on-chain-order-books-and-perpetual-dexs:8}
Les \omterm{def:m3:blockchains-for-traders:pow}{validateurs} peuvent voter la radiation d'un marché et le régler au prix qu'ils fixent selon la règle, et les paramètres du protocole (marges, caps, sources de l'\omterm{def:m3:blockchains-for-traders:contract}{oracle}) peuvent changer par la gouvernance.
Décentralisé signifie qu'aucune société unique ne décide, et non que rien ne puisse changer.
\end{solution}

\begin{solution}{pb:m3:on-chain-order-books-and-perpetual-dexs:1}
\textbf{1.} La moitié de la marge initiale de 33.3~\% à 3 fois~: 16.7~\%~; aux deux tiers de ce montant, 11.1~\% du notionnel, soit environ 1.11~million de dollars.
\textbf{2.} Le carnet était trop mince~: un ordre au marché de 10~millions de dollars aurait fait beaucoup bouger le prix, et il a été envoyé par tranches de 20~\% qui ne rétablissaient pas
la marge. \textbf{3.} 8.89~millions de dollars à 100~\%~; 38.89~millions de dollars à 400~\%. \textbf{4.} Il détient des positions longues de compensation dans d'autres
comptes ou marchés qui gagnent à la hausse, tandis que la perte sur la position vendeuse liquidée s'arrête à sa marge. \textbf{5.} Déplacer le cours au comptant sur
plusieurs plateformes externes à la fois, puisque l'\omterm{def:m3:blockchains-for-traders:contract}{oracle} est une médiane pondérée calculée sur elles~: les plateformes les moins coûteuses à déplacer fixent le coût de l'attaque.
\textbf{6.} $5/(5 - 0.111) \approx$ 1.02~million de dollars. \textbf{7.} 7.78~millions de dollars. \textbf{8.} Le short squeeze que l'attaquant peut se permettre dépend
de ce qu'il coûte de déplacer les sources de l'\omterm{def:m3:blockchains-for-traders:contract}{oracle}~; des marchés externes peu liquides rendent bon marché les grands short squeezes. \textbf{9.} Des paliers de marge plus élevés pour les grandes
positions, un effet de levier maximal plus faible, des plafonds globaux de positions ouvertes, les sources et pondérations de l'\omterm{def:m3:blockchains-for-traders:contract}{oracle}, et la taille des tranches de liquidation. \textbf{10.} Les
plafonds, paliers et règles de l'\omterm{def:m3:blockchains-for-traders:contract}{oracle} qui bornent sa queue de risque, et le processus de gouvernance qui peut les modifier. \textbf{11.} Toutes les positions sont réglées au
prix fixé par la règle, et les ordres sont annulés. \textbf{12.} Les positions vendeuses sur l'actif comprimé gagnent par rapport à un règlement au prix comprimé, les
positions longues de l'attaquant perdent~; le coffre est épargné par la poursuite du short squeeze. \textbf{13.} Les deux~: c'est une décision collective des détenteurs de \omterm{def:m3:blockchains-for-traders:key}{jetons} de gouvernance, et non d'une
société, mais elle montre que les règles peuvent changer pendant qu'une position est ouverte. \textbf{14.} Par son propre personnel~: relèvement des marges, restriction du marché,
radiation, ou règlement à un prix choisi selon ses conditions d'utilisation. \textbf{15.} Sous la forme d'une décote ou d'une limite sur les positions dans les marchés où de tels
votes sont plausibles (marchés peu liquides, récemment cotés, adossés à un coffre). \textbf{16.} Seulement avec des caps dimensionnés à la profondeur des marchés externes, ou pas du tout
pour les \omterm{def:m3:blockchains-for-traders:key}{jetons} dont le cours au comptant peut être déplacé à bas coût. \textbf{17.} Plus sûr~: la conservation on-chain, un état public, des règles qui s'appliquent à tous~; moins sûr~:
la dépendance à l'\omterm{def:m3:blockchains-for-traders:contract}{oracle}, des filets de sauvegarde publics pouvant être visés, l'intervention de la gouvernance, et le risque lié aux \omterm{def:m3:blockchains-for-traders:contract}{smart contracts}. \textbf{18.} Là où la \omterm{def:m3:on-chain-order-books-and-perpetual-dexs:priority}{règle de priorité dans le
bloc} protège ses cotations et où les frais conviennent à son volume, typiquement le \omterm{def:m3:on-chain-order-books-and-perpetual-dexs:perpdex}{carnet d'ordres on-chain} avec annulations en premier~; le pool à prix d'\omterm{def:m3:blockchains-for-traders:contract}{oracle} est une
contrepartie, et non un lieu pour faire de la tenue de marché. \textbf{19.} \emph{Résultat à retenir~:} une perte de $10\text{ millions} \times (s - 0.111)$ de dollars pour un
short squeeze $s$ (38.89~millions de dollars à 400~\%), et un cap d'environ 1.02~million de dollars pour la maintenir sous 5~millions de dollars à 500~\%. \textbf{20.} Le capital des déposants,
qui soutient chaque liquidation, et donc une cible pour quiconque peut faire échouer une liquidation.
\end{solution}

\begin{solution}{iq:m3:on-chain-order-books-and-perpetual-dexs:1}
Elle supprime l'exécution opportuniste des cotations périmées chaque fois que l'annulation du teneur de marché figure dans le même bloc que l'ordre du preneur, ce qui réduit la sélection adverse~;
les teneurs de marché peuvent coter plus serré et plus gros, et les preneurs doivent devancer le teneur de marché d'un bloc entier pour exécuter une cotation.
\iqlookfor{la sélection adverse passe des teneurs de marché aux preneurs de liquidité.}
\end{solution}

\begin{solution}{iq:m3:on-chain-order-books-and-perpetual-dexs:2}
Les messages sont des transactions signées dotées d'un nonce~; les accusés de réception arrivent bloc par bloc~; l'état des ordres se lit dans les blocs, et non dans un
flux privé~; les règles d'ordonnancement dans le bloc déterminent l'appariement~; le \omterm{def:m3:blockchains-for-traders:gas}{gas} ou les \omterm{def:m3:centralised-exchanges:ratelimit}{limites de débit} s'appliquent à chaque action~; il faut aussi rapprocher les données de l'état de la chaîne~;
et gérer les clés des portefeuilles de signature.
\iqlookfor{les blocs comme horloge, et les signatures comme session.}
\end{solution}

\begin{solution}{iq:m3:on-chain-order-books-and-perpetual-dexs:3}
La médiane ignore les valeurs aberrantes jusqu'à la moitié de leur poids, de sorte qu'une plateforme manipulée ou défaillante ne la déplace pas~; une moyenne pondérée par les volumes suit chaque
source au prorata et peut être déplacée par des volumes de \omterm{def:m3:spot-markets:wash}{wash trading}. La médiane peut sauter d'une source à l'autre et retarde quand le marché se déplace sur les
sources les plus pondérées.
\iqlookfor{robustesse face à la réactivité.}
\end{solution}

\begin{solution}{iq:m3:on-chain-order-books-and-perpetual-dexs:4}
Les dépôts dans un coffre ou un pool s'analysent comme un crédit et une exposition de queue~: une limite par plateforme, dimensionnée sur les pertes de stress (short squeezes, défaillances de l'\omterm{def:m3:blockchains-for-traders:contract}{oracle}), des périodes de blocage,
un suivi des positions ouvertes et des caps, et des propositions de gouvernance.
\iqlookfor{pertes de queue et liquidité du dépôt.}
\end{solution}

\begin{solution}{iq:m3:on-chain-order-books-and-perpetual-dexs:5}
Attaque~: négocier lorsque l'\omterm{def:m3:blockchains-for-traders:contract}{oracle} est périmé (latence), déplacer des sources de l'\omterm{def:m3:blockchains-for-traders:contract}{oracle} peu liquides, ou ouvrir des positions à un prix dont on sait qu'il va bientôt changer. Défense~:
des \omterm{def:m3:blockchains-for-traders:contract}{oracles} multi-sources robustes, des délais ou frais d'impact sur le prix, des \omterm{def:m3:on-chain-order-books-and-perpetual-dexs:oicap}{plafonds de positions ouvertes}, et des limites sur la taille des positions au regard de la profondeur des sources.
\iqlookfor{l'écart entre l'\omterm{def:m3:blockchains-for-traders:contract}{oracle} et le prix réel.}
\end{solution}

\begin{solution}{iq:m3:on-chain-order-books-and-perpetual-dexs:6}
Ingérer les blocs avec leurs transactions dans l'ordre proposé~; implémenter exactement la \omterm{def:m3:on-chain-order-books-and-perpetual-dexs:priority}{règle de priorité dans le bloc} de la plateforme et l'appariement~; rejouer à partir d'un instantané connu,
en vérifiant le carnet et les positions obtenus par rapport à l'état publié aux points de contrôle~; garder le rejoueur déterministe (arithmétique entière, tri
stable) et le versionner avec les mises à niveau du protocole.
\iqlookfor{règles exactes, instantanés et points de contrôle.}
\end{solution}
